\chapter{Procedure tipiche}
\label{ch:workflow_tipici}
\index{Workflow!tipico}\index{Procedure}\index{Casi d'uso}

\epigraph{``Per imparare a far bene una cosa nuova, conviene
ripeterla cinque volte. La prima volta a tentoni, la seconda
con cautela, la terza con scioltezza, la quarta con grazia,
la quinta con maestria.''}{Adattato da \textsc{Beatrice
Bizzarri}, \emph{Quaderno di didattica}, 1923}

\lettrine[lines=3]{Q}{uesto} capitolo raccoglie le sei procedure
ricorrenti dell'esperienza di un coordinatore d'orario. Ciascuna
\`e descritta come una sequenza di passi numerati, con la tab
da aprire, i bottoni da premere, l'effetto atteso. Sono il
ricettario operativo: per imparare a memoria una scuola di
trentadue classi in due ore. Le procedure sono indipendenti
l'una dall'altra: una scuola che apre l'orario ad agosto seguir\`a
le prime due; una scuola che subentra in corso d'anno con un
file Excel salter\`a alla terza.

\section{Procedura 1: una scuola da zero}
\label{sec:wf_da_zero}

Partenza: database vuoto, scuola da configurare per la prima
volta. Arrivo: orario settimanale completo visualizzato.

\subsection{Passo 1 -- configurare la griglia oraria}

Vai sulla tab \emph{Ore}. Verifica che i giorni lavorativi e
gli slot orari rispecchino la realt\`a della scuola. Il default
(lun--sab, 8:00--14:00, sei slot) va bene per la stragrande
maggioranza delle scuole italiane. Modifica solo se necessario.
Premi \emph{Salva}.

\paragraph{Attenzione.} Fissa la griglia ora, non dopo.
Cambiarla in seguito invalida le matrici di indisponibilit\`a
gi\`a inserite.

\subsection{Passo 2 -- caricare l'anagrafica}

Hai due strade.

\paragraph{Strada A: import da Excel.}
Vai sulla tab \emph{Import bulk}. Scarica i template dei sette
fogli (\texttt{teachers}, \texttt{subjects}, \texttt{classes},
\texttt{classrooms}, \texttt{curricula}, \texttt{students},
\texttt{groups}). Compila ciascun foglio. Reimporta nello stesso
ordine, modalit\`a \texttt{upsert}.

\paragraph{Strada B: inserimento manuale.}
Visita le tab \emph{Indirizzi}, \emph{Materie}, \emph{Docenti},
\emph{Classi}, \emph{Aule}, \emph{Studenti}, \emph{Gruppi} in
quest'ordine, e usa il bottone \emph{+ Nuovo} di ciascuna per
inserire le righe. Pi\`u lungo, ma adatto a scuole piccole o a
quando manca un foglio Excel di partenza.

\subsection{Passo 3 -- assegnare le cattedre}

Vai sulla tab \emph{Cattedre}. Per ogni classe vedi una card
con tutte le materie del suo indirizzo. Premi \emph{cambia}
accanto a ogni materia e scegli il docente. Per la prima
sessione conviene partire dai docenti pi\`u vincolati (esempio:
chi ha una sola classe possibile), e finire con quelli
\emph{jolly}. A configurazione completa, controlla con il
bottone \emph{Warnings cattedre} che nessun docente sia
\texttt{SFORA} o \texttt{SOTTO}.

\subsection{Passo 4 -- aggiungere indisponibilit\`a e vincoli}

Per ogni docente apri la scheda dalla tab \emph{Docenti} e
colora la sua matrice. Usa il rosso (\textsc{Hard}) per le
indisponibilit\`a vere; il giallo (\textsc{Soft}) per le
preferenze negative; il blu (\textsc{Preferred}) per i giorni
preferiti positivi.

Per i vincoli logici (es.\ ``Bianchi non lavora il sabato'',
``in 5A storia non al venerd\`{\i}'') vai sulla tab
\emph{Vincoli}, premi \emph{+ Nuovo vincolo DSL}, scrivi
l'espressione, \emph{Valida}, \emph{Salva}.

\begin{esempiobox}[title={Vincoli ricorrenti, in DSL}]
\begin{lstlisting}
# Bianchi non lavora il sabato
forall l in lessons where l.teacher == Bianchi:
    l.day != sab

# In 5A storia non al venerdi'
forall l in lessons where l.class == 5A and l.subject == Storia:
    l.day != ven

# La palestra e' occupata il martedi' mattina
forall l in lessons where l.classroom == Palestra:
    l.day != mar or l.hour > 11
\end{lstlisting}
\end{esempiobox}

\subsection{Passo 5 -- snapshot prima del calcolo}

Vai sul \emph{Dashboard}, card \emph{Import / Export DB}, premi
\emph{Salva snapshot}. Cos\`{\i} hai un ritorno garantito se la
prima pipeline produce risultati strani.

\subsection{Passo 6 -- pre-check}

Sulla tab \emph{Diagnostica} premi \emph{Hall pre-check}. Il
sistema verifica che le ore richieste siano formalmente
servibili dai docenti disponibili. Se il check fallisce, il
pannello indica il sottoinsieme di classi e materie in deficit:
non lanciare la pipeline, correggi prima.

\subsection{Passo 7 -- lanciare la pipeline}

Vai sulla tab \emph{Workflow}, premi \emph{Pipeline completa}.
Il pannello di log inferiore mostra le fasi una dopo l'altra
(Phase A, Phase B per giorno, metaeuristiche, aule). A
seconda della dimensione della scuola:

\begin{itemize}
\item \texttt{small} (8 classi): 30--90 secondi;
\item \texttt{medium} (15--20 classi): 2--5 minuti;
\item \texttt{huge} (40+ classi): 8--20 minuti.
\end{itemize}

Al termine il banner verde \emph{Pipeline completata} ti porta
direttamente alla tab \emph{Orario}.

\subsection{Passo 8 -- visualizzare e raffinare}

Sulla tab \emph{Orario}, naviga le quattro viste (per classe,
docente, aula, slot). Se trovi una lezione fuori posto, trascinala
in un altro slot: il sistema verifica in tempo reale i vincoli
Hard. Se tutto bene, esporta con il bottone \emph{Esporta XLSX}.

\section{Procedura 2: importare dati da Excel di un'altra scuola}
\label{sec:wf_import_excel}

Scuola che riceve da un altro istituto, o da un anno precedente,
i fogli Excel dell'anagrafica.

\paragraph{Premessa.} Non tutti i fogli arrivano nel formato del
nostro template. Conviene un piccolo lavoro di rinomina prima
dell'import.

\begin{enumerate}
\item Apri il template ufficiale dalla tab \emph{Import bulk}
  (bottone \emph{Scarica template} per \texttt{teachers}). Apri
  il foglio \emph{Istruzioni}: l\`{\i} sono elencati i nomi
  canonici e i sinonimi accettati (es.\ \texttt{nome},
  \texttt{teacher\_name}, \texttt{cognome\_nome} sono tutti
  accettati per il campo \texttt{name}).
\item Sul tuo Excel di partenza, rinomina le intestazioni
  delle colonne con i nomi canonici (o uno qualsiasi degli
  alias). Le colonne extra vengono ignorate.
\item Salva il file. Sulla tab \emph{Import bulk} scegli
  \texttt{teachers}, modalit\`a \texttt{upsert}, trascina il
  file. Premi \emph{Importa}.
\item Il pannello-report sotto il bottone elenca le righe
  importate e quelle scartate. Le scartate riportano l'indice
  e il motivo dell'errore (campo mancante, valore fuori
  range, eccetera).
\item Ripeti per gli altri sei fogli, nello stesso ordine.
\end{enumerate}

\begin{erroricomunibox}
\textbf{``Il file ha gli accenti rovinati.''} L'importer
accetta Excel xlsx (consigliato) o CSV. Se il CSV viene da
Excel italiano, salvalo in formato \emph{CSV UTF-8 (delimited)};
altrimenti la lettera \`e diventa una sequenza di simboli e i
record vengono scartati.

\textbf{``Le classi sono importate ma non hanno indirizzo.''}
Probabilmente nel file mancava la colonna \texttt{curriculum}
oppure il valore non corrispondeva a nessun indirizzo del
database. Sistema l'errore importando prima gli indirizzi e
poi le classi.
\end{erroricomunibox}

\section{Procedura 3: rilanciare la pipeline dopo un cambio di vincoli}
\label{sec:wf_relancia}

Scenario tipico: il collegio dei docenti chiede una modifica
sostanziale (es.\ ``aggiungiamo una compresenza in tutte le
prime''). La pipeline va rilanciata.

\begin{enumerate}
\item Snapshot della soluzione corrente, dal \emph{Dashboard},
  card \emph{Import / Export DB}, bottone \emph{Salva snapshot}.
  Lo snapshot serve a confrontare prima/dopo.
\item Aggiungi i nuovi vincoli (matrice, DSL, compresenze\dots).
\item Sulla tab \emph{Diagnostica} premi \emph{Hall pre-check}
  per verificare che la richiesta sia ancora formalmente
  fattibile.
\item Vai sulla tab \emph{Workflow}, premi \emph{Pipeline
  completa}. Se la prima pipeline aveva impiegato N minuti, la
  seconda impiegher\`a un tempo simile.
\item Confronta. Sulla tab \emph{Runs} vedrai la run nuova in
  alto. Cliccala per il dettaglio, confronta il SOFT score con
  la precedente.
\item Se la nuova \`e migliore: \`e fatta. Se peggiore: vai
  sullo snapshot precedente e premi \emph{Ripristina}.
\end{enumerate}

\section{Procedura 4: editare a mano una singola lezione}
\label{sec:wf_edit_manuale}

Scenario: l'orario \`e in produzione da una settimana, il
preside chiede di spostare l'inglese di 2B dal venerd\`{\i} al
gioved\`{\i}.

\begin{enumerate}
\item Sulla tab \emph{Orario}, scegli la vista \emph{Per
  classe}, seleziona 2B.
\item Trascina la cella di Inglese dal venerd\`{\i} al
  gioved\`{\i}. Mentre trascini, le celle candidate del
  gioved\`{\i} si colorano: verdi se accettabili, gialle se
  peggiorano il SOFT, rosse se violano un Hard.
\item Lascia la presa sopra uno slot verde. L'orario si
  aggiorna; il pannello \emph{Cost} in alto cambia.
\item Se l'aula non era compatibile col nuovo slot, si apre
  il modal \emph{Aula da scegliere}: seleziona una delle
  alternative verdi.
\item Esporta il nuovo orario (XLSX) per distribuirlo.
\end{enumerate}

\paragraph{Suggerimento.} La modifica manuale non aggiorna
automaticamente i vincoli salvati. Se la modifica deve
diventare permanente (\emph{l'inglese di 2B sta sempre al
gioved\`{\i}}), conviene aggiungere il vincolo corrispondente
nella tab \emph{Vincoli} cos\`{\i} la prossima pipeline lo
rispetter\`a senza dover ripetere il drag-and-drop.

\section{Procedura 5: gestire una supplenza giornaliera}
\label{sec:wf_supplenza}

Scenario quotidiano: tre docenti sono assenti, occorre
piazzare le supplenze per oggi.

\begin{enumerate}
\item Apri la tab \emph{Assenze e supplenze}.
\item Clicca l'intestazione del giorno corrente (esempio:
  \emph{Marted\`{\i}}). Si apre un modal con la lista dei
  docenti; spunta gli assenti, conferma.
\item Le celle dell'orario di quei docenti diventano rosse:
  classi scoperte.
\item Per ogni cella rossa, click sulla cella. Si apre un
  modal a due colonne: \emph{classi scoperte} a sinistra,
  \emph{docenti disponibili} a destra. I docenti disponibili
  sono ordinati per \texttt{graduatoria\_score}: i pi\`u
  indicati appaiono in cima.
\item Trascina un docente sulla classe scoperta; la cella
  diventa verde.
\item Se tutte le celle sono verdi, la giornata \`e coperta.
  Esporta il bollettino assenze con il bottone in alto a
  destra; verr\`a stampato e affisso.
\end{enumerate}

\begin{esempiobox}[title=Caso reale]
Liceo Galileo, marted\`{\i} ore 7:45. Sono assenti Bianchi
(matematica, 4 ore), Rossi (storia, 2 ore) e Verdi (inglese,
3 ore): nove celle scoperte. Il sistema propone in cima alle
liste tre supplenti coerenti per ciascuna materia, ordinati
per punteggio. La copertura si chiude in dodici minuti.
\end{esempiobox}

\section{Procedura 6: esportare e archiviare l'orario}
\label{sec:wf_export}

A fine anno, o quando si chiude una versione consolidata,
conviene archiviare tre cose:

\begin{enumerate}
\item L'\emph{orario in XLSX}. Vai sulla tab \emph{Orario},
  vista \emph{Per classe}, premi \emph{Esporta XLSX}. Otterrai
  un file con un foglio per classe, formattato come modulo
  classico delle scuole italiane. Ripeti per la vista \emph{Per
  docente} se serve.
\item Il \emph{database completo}. Vai sul \emph{Dashboard},
  card \emph{Import / Export DB}, premi \emph{Esporta DB
  completo}. Otterrai uno zip che contiene il file SQLite raw,
  i CSV per tabella e i metadati. \`E il backup ufficiale.
\item I \emph{vincoli}. Sulla stessa card \emph{Import / Export
  vincoli}, premi \emph{Esporta JSON} (oppure xlsx). Avrai un
  file con tutti i vincoli logici scritti dall'utente, comodi
  per il riuso in scuole simili.
\end{enumerate}

\paragraph{Nota di archiviazione.} Conviene salvare i tre file
in una cartella nominata \texttt{<nome\_scuola>\_<aa\_scolastico>}
sul drive condiviso della scuola. Anno per anno, l'archivio
diventa la base per studi storici (es.\ ``dove abbiamo
peggiorato/migliorato'') e per ripartire l'anno successivo con
una struttura nota.

\fleuron

\begin{biblioparvabox}
Le sei procedure di questo capitolo coprono il novanta per
cento dei casi d'uso reali. Le situazioni residue
(diagnostica avanzata, infeasibility profonda, branch-and-price
per scuole enormi) sono trattate nei capitoli specialistici:
\cref{ch:diagnostica_statistica} per la diagnostica statistica,
\cref{ch:tecniche_avanzate} per le tecniche
avanzate, \cref{ch:lessons_learned} per gli scenari
problematici incontrati nelle scuole pilota.
\end{biblioparvabox}
